\documentclass[11pt]{article}
% ===========================================================================
%  Shared preamble for the AP Calculus AB exam-preparation worksheets.
%  Input right after \documentclass. Files beginning with "_" are not compiled
%  on their own. Navy and gold, to match the course's lesson pages.
% ===========================================================================
\usepackage[letterpaper,top=2.3cm,bottom=2.2cm,left=1.7cm,right=1.7cm,headheight=28pt,footskip=30pt]{geometry}
\usepackage{amsmath,amssymb}
\usepackage[T1]{fontenc}
\usepackage[utf8]{inputenc}
\usepackage{lmodern}
\usepackage{xcolor}
\usepackage[most]{tcolorbox}
\usepackage{enumitem}
\usepackage{fancyhdr}
\usepackage{array}
\usepackage{tikz}
\usepackage{pgfplots}
\pgfplotsset{compat=1.18}

\definecolor{navy}{HTML}{161E54}
\definecolor{ink}{HTML}{27357E}
\definecolor{gold}{HTML}{B98A00}
\definecolor{goldbg}{HTML}{FFFDF3}
\definecolor{inkbg}{HTML}{F6F7FD}
\definecolor{rule}{HTML}{C9CFF0}
\definecolor{okgreen}{HTML}{1B7A44}
\definecolor{okbg}{HTML}{F1F8F3}

\pagestyle{fancy}
\fancyhf{}
\fancyhead[L]{\footnotesize Name: \underline{\hspace{4cm}}}
\fancyhead[C]{\textbf{\color{navy}Integration Academy}\\[-2pt]{\scriptsize Calculus AB \textperiodcentered{} AP\textsuperscript{\textregistered} Exam Preparation}}
\fancyfoot[L]{\tiny\color{black!55}\parbox[t]{0.78\textwidth}{AP\textsuperscript{\textregistered} is a registered trademark of the College Board, which was not involved in the production of, and does not endorse, this material. \textcopyright{} Integration Academy.}}
\fancyfoot[R]{\scriptsize\thepage}
\renewcommand{\headrulewidth}{1.2pt}
\renewcommand{\headrule}{\hbox to\headwidth{\color{gold}\leaders\hrule height \headrulewidth\hfill}}
\renewcommand{\footrulewidth}{0pt}
\setlength{\parindent}{0pt}
\setlength{\parskip}{4pt}

% Title block: \wstitle{code}{title}{unit line}
\newcommand{\wstitle}[3]{\begin{center}{\small\bfseries\color{gold}WORKSHEET #1}\\[3pt]{\LARGE\bfseries\color{navy}#2}\\[3pt]{\small\color{black!60}#3}\end{center}\vspace{2pt}}
% Section heading
\newcommand{\sech}[1]{\par\vspace{10pt}{\large\bfseries\color{navy}#1}\par\nobreak\vspace{2pt}{\color{rule}\hrule height0.8pt}\vspace{7pt}}
% A level tag printed beside a group of questions
\newcommand{\level}[1]{\par\vspace{6pt}{\footnotesize\bfseries\color{gold}\MakeUppercase{#1}}\par\nobreak\vspace{2pt}}

\newtcolorbox{ideabox}{colframe=ink,colback=inkbg,coltitle=white,colbacktitle=ink,fonttitle=\bfseries\small,title={Key ideas},arc=2pt,boxrule=0.8pt,left=8pt,right=8pt}
\newtcolorbox{example}[1]{colframe=rule,colback=white,coltitle=navy,colbacktitle=inkbg,fonttitle=\bfseries\small,title={#1},arc=2pt,boxrule=0.8pt,borderline west={3pt}{0pt}{ink},left=10pt,right=8pt}
\newtcolorbox{frbox}[1]{colframe=gold,colback=goldbg,coltitle=white,colbacktitle=gold,fonttitle=\bfseries\small,title={#1},arc=2pt,boxrule=0.8pt,breakable,left=8pt,right=8pt}
\newtcolorbox{solbox}{colframe=okgreen,colback=okbg,coltitle=white,colbacktitle=okgreen,fonttitle=\bfseries\small,title={Worked solutions},arc=2pt,boxrule=0.8pt,breakable,left=8pt,right=8pt}
\newtcolorbox{scorebox}[1]{colframe=navy,colback=white,coltitle=white,colbacktitle=navy,fonttitle=\bfseries\small,title={#1},arc=2pt,boxrule=0.8pt,breakable,left=8pt,right=8pt}

\newcommand{\soln}{\par\smallskip\textbf{\color{ink}Solution.}\ }
\newcommand{\result}[1]{\par\smallskip\textbf{#1}}
% blank working space of a given height
\newcommand{\work}[1]{\par\vspace{#1}}
% a point value in the margin of a scoring guide
\newcommand{\pts}[1]{\hfill{\footnotesize\bfseries\color{gold}[#1]}}

% Coordinate grid for graph questions: \begin{apgraph}{xmin}{xmax}{ymin}{ymax} ... \end{apgraph}
\newenvironment{apgraph}[4]{\begin{center}\begin{tikzpicture}\begin{axis}[
  axis lines=middle, xmin=#1, xmax=#2, ymin=#3, ymax=#4,
  grid=both, grid style={black!12}, xtick distance=1, ytick distance=1,
  ticklabel style={font=\scriptsize}, xlabel={$x$}, ylabel={$y$},
  width=10cm, height=7.2cm, samples=80, clip=false,
  every axis plot/.append style={line width=1.1pt, color=ink}]}{\end{axis}\end{tikzpicture}\end{center}}
\newcommand{\closeddot}[2]{\addplot[only marks, mark=*, mark size=2.2pt, color=ink] coordinates {(#1,#2)};}
\newcommand{\opendot}[2]{\addplot[only marks, mark=*, mark size=2.2pt, color=ink, mark options={fill=white}] coordinates {(#1,#2)};}

% Multiple-choice options: four across, or stacked when the options are long.
\newcommand{\opts}[4]{\par\vspace{3pt}\begin{tabular}{@{}p{0.235\linewidth}p{0.235\linewidth}p{0.235\linewidth}p{0.235\linewidth}@{}}(A)\ #1 & (B)\ #2 & (C)\ #3 & (D)\ #4\end{tabular}\par\vspace{9pt}}
\newcommand{\optsv}[4]{\par\vspace{3pt}\begin{tabular}{@{}l@{\ \ }p{0.9\linewidth}@{}}(A) & #1\\(B) & #2\\(C) & #3\\(D) & #4\end{tabular}\par\vspace{9pt}}

\fancyhead[R]{\footnotesize\color{black!60}Worksheet 1.7}
\begin{document}

\wstitle{1.7}{Where Graphs Break: Classifying Discontinuities}{Unit 1: Limits and Continuity}

\begin{ideabox}
\begin{itemize}[leftmargin=*,itemsep=2pt,topsep=2pt]
  \item \textbf{Removable:} the limit exists but $f(c)$ is undefined or different. A hole.
  \item \textbf{Jump:} both one-sided limits exist and are unequal.
  \item \textbf{Infinite:} a one-sided limit is $\pm\infty$. A vertical asymptote.
  \item For a rational function, factor first: a cancelled factor gives a hole; a surviving factor gives an asymptote.
\end{itemize}
\end{ideabox}

\sech{Worked examples}
\begin{example}{Example 1 \textperiodcentered{} A rational function}
Find and classify the discontinuities of $f(x)=\dfrac{x^2-4x}{x^2-16}$.\soln $f(x)=\dfrac{x(x-4)}{(x-4)(x+4)}$. At $x=4$ the factor cancels and $\displaystyle\lim_{x\to 4}f(x)=\dfrac48=\dfrac12$. At $x=-4$ the factor remains and the numerator tends to $-4\neq0$.\result{Removable at $x=4$; infinite at $x=-4$.}
\end{example}

\begin{example}{Example 2 \textperiodcentered{} A piecewise function}
Classify the discontinuity at $x=2$ of $g(x)=\begin{cases}2x-1,&x<2\\6,&x=2\\x+1,&x>2.\end{cases}$\soln $\displaystyle\lim_{x\to 2^-}g(x)=3$ and $\displaystyle\lim_{x\to 2^+}g(x)=3$, so the limit is 3; but $g(2)=6$.\result{Removable.}
\end{example}

\begin{example}{Example 3 \textperiodcentered{} An absolute value}
Classify the discontinuity at $x=3$ of $h(x)=\dfrac{x^2-9}{|x-3|}$.\soln For $x>3$: $h(x)=x+3\to6$. For $x<3$: $h(x)=-(x+3)\to-6$.\result{Jump.}
\end{example}

\newpage
\sech{Practice}

\level{Core}
\begin{enumerate}[leftmargin=*,itemsep=2pt]
  \item Classify the discontinuity of $f(x)=\dfrac{x^2-1}{x-1}$ at $x=1$.\work{2.3cm}
  \item Classify the discontinuity of $f(x)=\dfrac{3}{x-2}$ at $x=2$.\work{2.3cm}
  \item Classify the discontinuity at $x=0$ of $f(x)=\begin{cases}x,&x<0\\x+2,&x\ge0.\end{cases}$\work{2.3cm}
  \item Classify the discontinuity of $f(x)=\dfrac{x-5}{(x-5)^2}$ at $x=5$.\work{2.3cm}
\end{enumerate}

\level{Exam level}
\begin{enumerate}[leftmargin=*,itemsep=2pt,start=5]
  \item Find and classify all discontinuities of $f(x)=\dfrac{x^2+x-2}{x^2-x-6}$.\work{2.8cm}
  \item How many vertical asymptotes does the graph of $f(x)=\dfrac{x^2-9}{x^3-9x}$ have? Justify.\work{2.8cm}
  \item For what value of $a$ does $f(x)=\dfrac{x^2+ax-8}{x-4}$ have a removable discontinuity at $x=4$? Give the coordinates of the hole.\work{2.8cm}
  \item Classify the discontinuity of $f(x)=\dfrac{1-\cos x}{x}$ at $x=0$.\work{2.8cm}
\end{enumerate}

\begin{frbox}{Question 9 \textperiodcentered{} Free-response practice}
Let $g(x)=\begin{cases}\dfrac{x^2-1}{x-1},&x<1\\[2mm]k,&x=1\\[1mm]3-x,&x>1\end{cases}$ where $k$ is a constant.
\begin{enumerate}[label=(\alph*),leftmargin=*,itemsep=2pt]
  \item Find $\displaystyle\lim_{x\to 1^-}g(x)$ and $\displaystyle\lim_{x\to 1^+}g(x)$.\work{2.6cm}
  \item If $k=5$, what type of discontinuity does $g$ have at $x=1$? Give a reason.\work{2.6cm}
  \item For what value of $k$ is $g$ continuous at $x=1$?\work{2.6cm}
\end{enumerate}
\end{frbox}

\newpage
\begin{solbox}
\begin{enumerate}[leftmargin=*,itemsep=4pt]
  \item Removable: $f(x)=x+1$ for $x\neq1$, limit 2, $f(1)$ undefined.
  \item Infinite: numerator 3, denominator tends to 0.
  \item Jump: left limit 0, right limit 2.
  \item Infinite: $f(x)=\dfrac{1}{x-5}$ for $x\neq5$; one factor remains.
  \item $f(x)=\dfrac{(x+2)(x-1)}{(x-3)(x+2)}$. Removable at $x=-2$ (limit $\dfrac{-3}{-5}=\dfrac35$); infinite at $x=3$.
  \item One. $f(x)=\dfrac{(x-3)(x+3)}{x(x-3)(x+3)}=\dfrac1x$ for $x\neq\pm3$: holes at $x=\pm3$, an asymptote only at $x=0$.
  \item $16+4a-8=0$ gives $a=-2$. Then $x^2-2x-8=(x-4)(x+2)$ and the hole is at $(4,6)$.
  \item Removable: $f(0)$ is undefined but $\displaystyle\lim_{x\to 0}\dfrac{1-\cos x}{x}=0$.
  \item \begin{enumerate}[label=(\alph*),leftmargin=*,itemsep=3pt]
    \item For $x<1$, $g(x)=x+1$, so $\displaystyle\lim_{x\to 1^-}g(x)=2$. $\displaystyle\lim_{x\to 1^+}g(x)=3-1=2$.
    \item Removable, because $\displaystyle\lim_{x\to 1}g(x)=2$ exists but $g(1)=5\neq2$.
    \item $k=2$, so that $g(1)$ equals the limit.
  \end{enumerate}
\end{enumerate}
\end{solbox}

\end{document}
